% Page 046

\seite{46}

ein Knabe widmet sich dem Maurergewerbe; die Mäd\ohb
chen widmen sich der Hauswirthschaft. Aufgenommen wur-
den 8 Kinder, und zwar 5 Knaben und 3 Mädchen.
Alle Kinder sind katholisch; ein Knabe ist nicht voll\ohb
sinnig.

\pend
\abschnitt{Spaziergang}
\pstart

Am 30.\ August unternahmen die vier Schulen der Pfarrei\ob
Seffern gemeinschaftlich einen Spaziergang nach Schönecken.

\pend
\abschnitt{Ernte}
\pstart

An Halmfrüchten war die diesjährige Ernte zufriedenstellend.\ob
An Grummet war sie gering wegen der vielen Regen\ohb
fälle des Sommers. Die Witterung war eine nur im\ob
ganzen günstige, und fast nur Kälte und Gewitter machte\unsicher\ob
nothwendig. Obst und Kartoffeln ergaben\ob
einen mittleren Ertrag.

\pend
\abschnitt{Gesehen 29.11.09}
\pstart


\pend
\abschnitt{Revision}
\pstart

Die hiesige Schule wurde am 29.\ November 1909 durch\ob
Herrn Kreisschulinspektor Schulrath Lentz\unsicher{} aus Bitburg revidiert.

\pend
\abschnitt{Joh. Schmitt, Pfr}
\pstart


\pend
\begin{verse}
14/3.10.
\end{verse}
\pstart

\pend
\abschnitt{Osterprüfung}
\pstart

Die Osterprüfung wurde am 14.\ März durch den Ortsschulinspek\ohb
tor Herrn Pfarrer Schmitt aus Seffern abgehalten.

\pend
\jahresueberschrift{Das Schuljahr 1910-11}
\pstart

Entlassen wurden in diesem Jahre 5 Knaben. Dieselben widmen\ob
sich der Landwirthschaft. Aufgenommen wurden 3 Mädchen und\ob
1 Knabe. Die Schülerzahl beträgt 80, davon sind 39 Mädchen und\ob
41 Knaben. Alle Kinder sind katholisch; ein Knabe ist nicht vollsinnig.
\pend
